\section{Discussion}

MouseBrainBench addresses a narrow but consequential problem: computational
neuroscience increasingly produces models that predict neural measurements,
while the language used to describe those models often extends to biological
mechanism or digital-twin fidelity. The experiments show that these
interpretations can be separated operationally. Known-truth systems expose
which evidential component is absent, Allen preserves a strong reproducibility
result while rejecting topology and direction, Sensorium preserves predictive
utility while rejecting causal interpretation, and MICRONS admits a local
structure--function association after matched controls and replication.

The negative Allen result is scientifically useful. A less restrictive
benchmark might report cross-mouse correlation of 0.891 and split-half
correlation of 0.990 as validation of a model. MouseBrainBench instead asks
whether the Allen-informed topology performs better than alternatives and
whether a directed signature can be recovered. It does not. This outcome does
not contradict the biological organization of the visual system. It identifies
an observational limitation: repeatable population dynamics do not uniquely
determine the graph that generated them.

Sensorium illustrates a related distinction in modern Artificial Intelligence.
Static stimulus and context models improve prediction over mean and scrambled
controls, while temporal models and a compact neural network improve predictions
for Dynamic Sensorium. These are valid functional results. They are not evidence
that the architecture's internal variables correspond to synapses, cell types,
or causal operations in the animal. Treating predictive and mechanistic scores
as separate outputs allows the same model to be accepted for one purpose and
rejected for another without contradiction.

The bounded official Sensorium experiment also clarifies what software
integration demonstrates. Executing an official loader and architecture verifies
interoperability, but a restricted local training budget is not comparable with
a published competition result. The very low bounded correlations are therefore
neither evidence against the official architecture nor a competitive baseline.
They document that the integration path works and that a full-budget comparison
remains absent. This prevents an engineering smoke test from becoming an
inflated empirical claim.

The MICRONS result is the strongest positive outcome, but its novelty must be
stated precisely. Ding et al. already demonstrated like-to-like connectivity
with a more extensive biological analysis, including feature, receptive-field,
and axon--dendrite geometry controls. MouseBrainBench does not improve that
biological finding. Instead, the MICRONS experiment tests whether a compact,
predefined audit can recover a compatible local association and preserve its
proper scope. Replication across two non-overlapping windows and positive
unit-cluster bootstrap intervals provide external validity for the audit
pipeline.

The effect is not causal. Readout-location similarity can covary with cortical
area, layer, cell class, developmental organization, or unmeasured morphology.
Distance and degree matching address two major confounds but do not exhaust the
causal graph. In particular, the analysis does not compare axon--dendrite
opportunity at the resolution used by Ding et al., manipulate individual
connections, or demonstrate that the observed edges are necessary for the
functional descriptor. The permitted interpretation is therefore a replicated
local observational association.

The current MIS should also not be treated as a universal scalar definition of
mechanistic explanation. Its principal contribution is structural: evidence
blocks are non-interchangeable and passage criteria are explicit. The particular
criteria and thresholds remain claim-dependent. A molecular mechanism, a
neural-mass model, and an individual-level behavioural twin would require
different blocks. External preregistration and application by independent groups
are needed before MIS can be regarded as a community standard.

Conjunctive rules have an intentional cost. They reduce false promotion of
partial evidence but can reject scientifically interesting models when one
required measurement is unavailable. MouseBrainBench handles this by retaining
block scores and partial-positive labels. Failure of the global decision means
that the target claim is unsupported under the declared evidence; it does not
mean that the model, data, or partial result is useless. This semantic discipline
is important if the framework is to encourage stronger experiments rather than
merely produce negative labels.

The sensitivity analysis indicates that the main qualitative conclusions are
not artifacts of one exact threshold. Allen remains negative under 25\%
relaxation and tightening. Static Sensorium retains partial reliability and
topographic evidence for six of nine threshold combinations, failing when the
reliability threshold reaches 0.70. Dynamic Sensorium gains over mean response
are more stable than gains over temporal SVD, supporting predictive utility but
not superiority of a specific neural architecture. These patterns make the
limitations visible instead of hiding them in a single aggregate score.

Several aspects constrain generalization. MICRONS data originate from one mouse
and a restricted visual-cortical volume. Its three cohorts are independent
windows of units, not independent animals. Allen analyses use a controlled
subset of available sessions. Dynamic Sensorium contains five current and five
legacy mice in the local package. The study contains no intervention, no
behavioural closed loop, and no longitudinal calibration of an individual
animal. Consequently, it cannot support whole-brain, behavioural, causal, or
consciousness claims.

For PLOS Computational Biology, the value of the work depends on the framework
being useful beyond this exact collection of analyses. The software interfaces
are therefore organized around claims, evidence blocks, graphs, predictions,
perturbations, and provenance rather than around hard-coded dataset names.
External BMTK, SONATA, or The Virtual Brain simulations can be added as candidate
models without reimplementing their numerical engines. New datasets can add
evidence blocks while retaining the same publication-freeze contract.

The next scientific step is not to increase model complexity indiscriminately.
It is to add evidence unavailable here: independent-animal structure--function
replication, perturbations that distinguish competing mechanisms, and
full-budget predictive baselines. Those additions would test whether the audit
can move from observational association towards intervention-supported
mechanism. Until then, MouseBrainBench should be understood as a reproducible
claim-control framework for partial mouse-brain models, not as the mouse-brain
model itself.
